\subsection{Comparison maps and elementary obstructions}

Throughout this subsection, $B\mathcal C$ denotes the realization of the
ordinary nerve. Products in a monoid are written in composition order:
$ab$ means that $b$ is applied first. Our chain formulas instead list
arrows chronologically. These conventions are converted explicitly in
the certificate programs.

\begin{lemma}[Proved: deleting an object retract]
\label{CP2:A-retract-deletion}
Suppose $x\ne y$ are objects of a finite category $C$, and there are
arrows $r:x\to y$ and $s:y\to x$ with $sr=1_x$. Let $D$ be the full
subcategory obtained by deleting $x$. The inclusion $j:D\to C$ induces
a homotopy equivalence. It removes at least three arrows, including
identities. If $C$ is connected with $k$ objects and $n$ arrows, the
excess $n-(2k-1)$ drops by at least one.
\end{lemma}
\begin{proof}
For $c\in D$, the comma category $(j\downarrow c)$ has terminal object
$(c,1_c)$. In $(j\downarrow x)$ choose $(y,s)$. An arbitrary object is
$(d,f:d\to x)$, and the arrow $rf:d\to y$ gives a morphism
$(d,f)\to(y,s)$ because $srf=f$. For a comma morphism
$h:(d,f)\to(d',f')$ we have $f'h=f$, so $(rf')h=rf$. These arrows
are therefore the components of a natural transformation from the
identity functor to the constant functor at $(y,s)$. They contract the
classifying space of this comma category. Quillen's Theorem~A,
\S1 of~\cite{CP2-Quillen1973}, applies to the specified inclusion.
The arrows $1_x,r,s$ are distinct and all disappear. Connectedness
follows either from the homotopy equivalence or directly from the
retract paths. Subtracting $2(k-1)-1$ from the new arrow count proves
the excess assertion.
\end{proof}

\begin{remark*}
It would be incorrect to define a functor deleting $x$ by conjugating
arrows with $r,s$ without checking identities: that proposed map sends
$1_x$ to $rs$, which need not be $1_y$. The comma-category proof avoids
this error. The statement concerns homotopy equivalence of classifying
spaces, not equivalence of categories.
\end{remark*}

\begin{lemma}[Proved: two blocks and their cup products]
\label{CP2:A-two-block}
Let $C$ be a finite category whose objects are partitioned into full
subcategories $D,E$, with no arrows from $E$ to $D$. There is a
double-bar space $T$ and a natural homotopy-pushout description
\[
 BC\simeq BD\mathop{\cup}\limits_T^h BE.
\]
If two positive-degree cohomology classes over a field both restrict
to zero on $BD$ and $BE$, their product is zero. In particular, if
$H^2(BD;\mathbb Q)=H^2(BE;\mathbb Q)=0$, all rational degree-two
products on $BC$ vanish.
\end{lemma}
\begin{proof}
Use the bisimplicial set whose $(p,q)$-simplices consist of a string
of $p$ arrows in $D$, an arrow from its last object to the first
object of a string of $q$ arrows in $E$, and that latter string.
Composition supplies the two actions and all faces; identities supply
degeneracies. Its iterated realization is $T$. Projection to the two
strings gives maps $T\to BD$ and $T\to BE$.

A simplex of $NC$ either belongs to one block or crosses exactly once
from $D$ to $E$. For a mixed simplex, sum the barycentric coordinates
on the $D$ vertices. Away from the two endpoints, division by this
sum and its complementary sum gives the two simplex coordinates and
one interval coordinate. At an endpoint the irrelevant block is
forgotten. These descriptions commute with faces and degeneracies.
Realizing them identifies $BC$ with the double-mapping-cylinder model,
up to the usual subdivision between diagonal and iterated realization
of the bisimplicial set. This proves the natural comparison asserted.

In that model put $A=BD\sqcup BE$. The inclusion is a cofibration, and
collapsing all of $A$ to one point gives $\Sigma(T_+)$. A class whose
restriction to $A$ is zero lifts to relative cohomology by the exact
sequence of the pair. All positive-degree products in a suspension
are zero: cover the suspension by its two contractible cones, lift
the two factors relative to the respective cones, and use the
relative cup product into the cohomology relative to their union,
which is zero. Applying the quotient identification and then the map
from relative to absolute cohomology proves the assertion.
\end{proof}

\begin{corollary}[Proved: a fixed incoming arrow]
\label{CP2:A-fixed-arrow}
Let $E$ consist of one object $v$ with corner monoid $M$. Suppose
$H^2(BD;\mathbb Q)=0$, and suppose some $f:d\to v$ satisfies
$mf=f$ for every $m\in M$. Then every rational degree-two product
on $BC$ is zero. The dual assertion holds for a fixed outgoing arrow.
\end{corollary}
\begin{proof}
The fixed arrow defines a simplicial section $BM\to T$ of the
projection $T\to BM$: use the object $d$ and arrow $f$ in the
$D$ and correspondence coordinates and retain the complete $M$
string. The equations $mf=f$ verify the face identities. Pullback
$H^2(BM;\mathbb Q)\to H^2(T;\mathbb Q)$ is therefore injective.
Mayer--Vietoris and the vanishing on $D$ imply that every degree-two
class of the pushout restricts to zero on both endpoint spaces.
Lemma~\ref{CP2:A-two-block} applies.
\end{proof}

\begin{lemma}[Proved: elementary facts about the minimal ideal]
\label{CP2:A-small-minimal-ideal}
Let $M$ be a finite monoid, $I$ a minimal nonempty two-sided ideal,
and $e\in I$ an idempotent. Put $G=eMe$. Then $G$ is a group.
If $Me=G$ or $eM=G$, the specified inclusion of the group corner in
the Karoubi envelope gives a homotopy equivalence $BG\simeq BM$.
If neither equality holds and $1\notin I$, then $|M|\ge5$.
At order five this last case is uniquely the identity-adjoined
$2\times2$ rectangular band.
\end{lemma}
\begin{proof}
Every finite ideal contains an idempotent: among the eventually
periodic powers of one of its elements choose a sufficiently large
multiple of the period. For $a\in eMe$, minimality gives $MaM=I$,
so $e=uav$ for suitable $u,v\in M$ and
$e=(eue)a(eve)$. In a finite monoid a product equal to the identity
forces its factors to be units: a one-sided inverse makes a finite
translation surjective and hence injective, giving the other
inverse. Applying this fact twice inside $eMe$ proves that it is a
group.

Let $K$ be the full two-object subcategory of $\operatorname{Kar}(M)$
on $1,e$ (one object when $e=1$). The object $e$ is a retract of $1$,
so Lemma~\ref{CP2:A-retract-deletion} identifies $BK$ with $BM$.
If $Me=G$, include the object $e$ into $K$. The comma category over
$1$ has objects the elements of $Me=G$ and arrows given by the free
transitive right $G$ action; there is exactly one arrow between
any ordered pair of objects. The comma category over $e$ has the
same description. Both are contractible, so Theorem~A gives
$BG\simeq BK$. If $eM=G$, use the dual comma categories.

Suppose neither equality holds. Choose $a\in Me\setminus G$.
Since $ea\in G$, replace $a$ by $a(ea)^{-1}$, obtaining
$ae=a$ and $ea=e$. Similarly choose $b\in eM\setminus G$ with
$eb=b$ and $be=e$. The four elements $e,a,b,ab$ are distinct:
the pairs of their left and right products with $e$ are respectively
$(e,e),(e,a),(b,e),(b,a)$. They lie in $I$, while $1$ lies outside.
Any element of $G$ other than $e$ is different from these four, since
its two products with $e$ both equal itself. Thus at order five
$G=\{e\}$ and $I=\{e,a,b,ab\}$. Now $ba\in eMe$ equals $e$,
and the relations just displayed determine the rectangular-band
table on these four elements. The remaining element is the
identity. Finally, if $1\in I$, then $I=M$ and the preceding
unit argument, with $e=1$, says that $M$ itself is a group.
\end{proof}

\begin{lemma}[Proved: rectangular bands with an identity]
\label{CP2:A-band-spheres}
For nonempty finite sets $I,J$ of sizes $r,c$, let
$M=\{1\}\sqcup\{x_{ij}:i\in I,j\in J\}$ with
$x_{ij}x_{kl}=x_{il}$. Then
\[
 BM\simeq\bigvee^{(r-1)(c-1)}S^2,
\]
where a wedge of zero spheres means a point.
\end{lemma}
\begin{proof}
Choose distinguished indices $0$ and put $R=\mathbb ZM$, $e=x_{00}$,
$P=Re$, $Q=R(1-e)$ and $D=R/\mathbb Z[I\times J]$. These are left
$R$-modules; $P,Q$ are projective because their defining elements
are orthogonal idempotents. The quotient $D$ is infinite cyclic as
an abelian group, with all band elements acting by zero.

The augmentation $P\to\mathbb Z$ has kernel with basis
$x_{i0}-x_{00}$ for $i\ne0$. Each band element annihilates these
differences, so this kernel is $D^{r-1}$ as an $R$-module.
The module $Q$ has the abelian basis
\[
 1-e,\qquad x_{ij}-x_{i0}\quad(i\in I,\ j\ne0).
\]
The map $Q\to D$ induced by the quotient has kernel the direct sum,
over $j\ne0$, of copies of $P$, with the map on the $j$th copy
$x_{i0}\mapsto x_{ij}-x_{i0}$. These maps are $R$-linear by the
rectangular-band multiplication rule. They are injective, and
their displayed bases prove exactness over the integers. Splicing
gives the complete projective resolution
\[
0\longrightarrow P^{(r-1)(c-1)}\longrightarrow Q^{r-1}
 \longrightarrow P\longrightarrow\mathbb Z\longrightarrow0.
\]
Tensor with the trivial right $R$-module $\mathbb Z$. The terms
$P,Q$ become respectively $\mathbb Z,0$, because the augmentations
of $e,1-e$ are $1,0$. Thus integral homology is $\mathbb Z$ in
degree zero, $\mathbb Z^{(r-1)(c-1)}$ in degree two, and zero
in every other positive degree. The universal group is trivial,
since every nonidentity element is idempotent. Hurewicz supplies a
map from the stated wedge inducing an isomorphism on $H_2$.
It is an integral homology equivalence between simply connected
CW spaces, hence a homotopy equivalence.
\end{proof}

\begin{corollary}[Proved: the order-five homotopy assertion in R7]
\label{CP2:A-five-types}
Every finite monoid of order at most five has a classifying space
homotopy equivalent to that of a finite group, except for the
five-element rectangular-band monoid, whose classifying space is
$S^2$. In particular, every corner of order at most four is
rationally acyclic, and every corner of order at most five has
zero third rational homology.
\end{corollary}
\begin{proof}
Apply Lemmas~\ref{CP2:A-small-minimal-ideal} and
\ref{CP2:A-band-spheres}. A finite group has zero positive rational
homology, since averaging splits its trivial rational module.
No finite enumeration is needed for this homotopy classification;
the separate census below verifies the table counts and identities.
\end{proof}

\begin{lemma}[Proved: exact projectivity test]
\label{CP2:A-projectivity-test}
For the preceding $M,e,G$, the trivial left $\mathbb QM$-module is
projective if and only if $Me=G$. The trivial right module is
projective if and only if $eM=G$.
\end{lemma}
\begin{proof}
Projectivity on the left is equivalent to a left-module splitting
of the augmentation $\mathbb QM\to\mathbb Q$: it suffices because
$\mathbb QM$ is free, and it is necessary by the defining lifting
property. Such a splitting is an invariant vector of augmentation
one. If $Me=G$, left multiplication by $m\in M$ on $G$ is group
translation by $me\in G$. The vector
$|G|^{-1}\sum_{g\in G}g$ supplies the splitting.

For the converse, first note that $eM$ is a minimal nonempty right
ideal. If $x\in eM$, write $e=uxv$ by minimality of the two-sided
ideal. Since $ex=x$, we have $e=(eue)xv$, so the inverse of
$eue\in G$ gives $e\in xM$. Thus $xM=eM$. The same argument
applies to every idempotent in $I$.
If $Me\ne G$, choose the normalized idempotent $a$ from the proof
of Lemma~\ref{CP2:A-small-minimal-ideal}. The minimal right ideals
$eM$ and $aM$ are disjoint: if they met they would be equal, and
then $a\in eM\cap Me=eMe=G$, a contradiction.
An invariant vector fixed by left multiplication by $e$ is
supported in $eM$, while one fixed by $a$ is supported in $aM$.
Hence an invariant vector must be zero and cannot have augmentation
one. This proves the left assertion. Opposition proves the right.
\end{proof}

\begin{lemma}[Proved: semisimple diagonal and support height]
\label{CP2:A-semisimple-height}
Suppose $C$ has acyclic support, and every corner algebra over
$\mathbb Q$ is semisimple. If every strict chain of objects has at
most $h$ arrows, then $H_i(BC;\mathbb Q)=0$ for $i>h$.
For a two-object category of acyclic support with at most five
arrows, $H^2(BC;\mathbb Q)=0$.
\end{lemma}
\begin{proof}
Let $A=\mathbb QC$, let $B$ be the product of the corner algebras,
and let $J$ be the span of all strict arrows. Acyclic support makes
$J$ an ideal and $A=B\oplus J$ an algebra splitting. The normalized
relative bar resolution of the constant left module $L$ has terms
$A\otimes_B J^{\otimes_B p}\otimes_B L$. Its standard insertion-of-
identity contraction proves exactness before augmentation. Because
$B$ is semisimple every $B$-module is projective, so each displayed
term is projective over $A$. A tensor in $J^{\otimes_B p}$ requires
a strict composable chain of $p$ objects changes; the tensor is
zero for $p>h$. Tensoring with the constant right module computes
the usual categorical nerve homology, giving the asserted bound.

For the last statement, write the two corner orders as $a,b$ and
the nonempty cross-hom size as $s$. If $s=1$, the double bar is
$BM\times BN$ and the collage is the join $BM*BN$. Both factors
are connected and have zero rational first homology (their
universal groups are finite), so the join has zero rational second
homology. If $s\ge2$, then $a+b\le3$, so both corner algebras
are semisimple: an order-two monoid is either $C_2$ or
$\{1,z:z^2=z\}$, both having rational algebra $\mathbb Q\times
\mathbb Q$. Apply the height-one case just proved.
\end{proof}

\begin{lemma}[Proved: singleton cross-homs]
\label{CP2:A-singleton-support}
Let $C$ be connected, have acyclic support and at least two objects,
and have one arrow in each nonempty cross-hom. If $P$ is the
support poset, the functor $q:C\to P$ induces an isomorphism on
rational second cohomology. This is the pullback of an actual
functor and therefore respects products.
\end{lemma}
\begin{proof}
Filter normalized rational chains by the number $p$ of strict
object changes. The internal complex for a chain
$x_0<\cdots<x_p$ is the tensor product of the corner bar complexes:
all intervening hom-sets are singleton trivial bimodules. Its
internal-degree-zero homology is $\mathbb Q$, and all internal
degree-one homology is zero, since every finite monoid has finite
universal group. Thus the bottom row of the resulting first-quadrant
spectral sequence is the chain complex of $BP$, and the next row
is zero.

In internal degree two, an edge contributes the direct sum of the
second homology of its two corners. Its horizontal boundary maps
each summand, with its boundary sign, onto that same corner
summand in horizontal degree zero. Every object is incident to
an edge by connectedness, so $E^2_{0,2}=0$. The only possible
survivor of total degree two is $E^2_{2,0}=H_2(BP;\mathbb Q)$.
No later differential enters it, and its only possible outgoing
differential lands in the zero internal-degree-one row.
The filtration has finitely many horizontal columns, so it
converges in each degree. The map induced by $q$ is the identity
on the bottom row, proving the homology and then cohomology
isomorphism. Naturality gives the assertion about products.
\end{proof}

\begin{lemma}[Proved: the support graph budget]
\label{CP2:A-cycle-budget}
Let $P$ be a finite connected poset, with comparability graph $G$
and cycle rank $R=|E(G)|-|V(G)|+1$.
If $R\le6$, then $H_4(BP;\mathbb Q)=0$.
If $R\le3$ and $H_1(BP;\mathbb F_2)=0$, then $G$ is chordal.
In the latter range its nontrivial clique blocks are triangles
in a clique tree, with at most three triangle additions, or one
complete four-vertex block.
\end{lemma}
\begin{proof}
The order complex is the clique complex of $G$. A four-simplex
requires a complete graph on five vertices, of cycle rank six.
If one occurs and $R\le6$, contracting it leaves a tree; all
remaining edges are bridges. There is exactly one four-simplex,
whose simplicial boundary is nonzero, and there is no
five-simplex. If none occurs, there are no four-chains at all.
Either case gives the first assertion.

For the second, triangle boundaries span the graph cycle space
over $\mathbb F_2$. Choose a basis of at most three triangle
boundaries. An induced cycle of length at least four has at most
one edge in any triangle: two would force a chord. No sum of
three basis triangles could contain all its edges, a
contradiction. Thus there is no induced cycle of length at least
four. Chordality gives the clique-tree decomposition, which can
also be read directly here by successive deletion of a simplicial
vertex. A four-clique uses all three cycles. Otherwise every
nontrivial maximal clique is a triangle, and adding a triangle
along an edge or a vertex adds one cycle. This proves the stated
small-range list.
\end{proof}

\subsection{A complete independent monoid census through order seven}

\begin{theorem}[Proved by complete finite computation]
\label{CP2:A-monoid-census}
The complete numbers of ordinary isomorphism classes of monoids of
orders one through seven are
\[
1,\ 2,\ 7,\ 35,\ 228,\ 2237,\ 31559.
\]
Up to isomorphism and anti-isomorphism the corresponding numbers are
\[1,2,6,27,156,1373,17730.\]
Our ordinary catalogs agree exactly,
after relabelling, with every supplied RESIDUE table in orders
one through five. These are counts of monoids with identity,
not counts of all semigroups or literal oriented table lists.
\end{theorem}
\begin{proof}
Our program \path{code/monoid_census.cpp} fixes the identity label
at zero and treats the other $(n-1)^2$ entries as unknown.
For each triple $(a,b,c)$ whose inner products are known it
enforces $(ab)c=a(bc)$, either rejecting an unequal pair of
known values or forcing one unknown value from the other.
Every such step is a necessary consequence of associativity.
It then takes the first remaining row-major unknown entry and
branches over every element of the monoid. Induction on the
search tree proves that every identity-zero associative table
occurs, and occurs once. The seven order-seven jobs partition
the search by the value of $1\cdot1$ and cover all seven
values disjointly.

At a complete table, intrinsic element invariants partition
possible relabellings. They include principal-ideal sizes,
translation fixed-point counts, commuting-element counts and
the power index and period. An isomorphism preserves these
invariants. We sort their values, enumerate every permutation
within an equal-invariant block, and retain the least resulting
word. Thus we remove exactly ordinary isomorphism duplicates;
opposition is not used in this enumeration.

The independent checker \path{code/check_monoid_census.cpp}
does not use this invariant partition. For every retained
table it enumerates all identity-fixing permutations, verifies
all associativity equations, checks that no two representatives
are isomorphic, computes the automorphism order, and verifies
that the recorded multiplicity is
$(n-1)!/|\operatorname{Aut}(M)|$. It also computes opposition
classes from the full permutation orbits. Finally
\path{code/verify_census.py} independently puts the supplied
RESIDUE tables into unrestricted lexicographic normal form and
compares the complete sets in each order at most five.

The labelled totals, with identity fixed, are
\[
1,\ 2,\ 11,\ 156,\ 4122,\ 208672,\ 18507440.
\]
The independent checker verifies $11,338,875$ associativity
equations in the representative tables. Every check passes;
the exact tables, automorphism-orbit checks, run counts,
source hashes and output hashes are supplied in
\path{outputs/monoid_census/}. Their ordinary and opposition
counts agree with the opened OEIS records
\cite{OEIS_A058129,OEIS_A058133} and the published distinction
between these two equivalence relations. The algorithmic
exhaustion proof above, not agreement of counts alone,
establishes completeness.
\end{proof}

\begin{remark*}[Catalog scope]
The number $1032$ used in the inherited order-eight search is
a restricted possible-positive-$H^2$ oriented cover, not a
complete ordinary order-eight monoid catalog. The published
ordinary count in order eight is $1,668,997$. Any use of that
restricted cover therefore requires its specific homological
reduction and an exhaustive proof within the reduced family.
The new order-one-through-seven catalogs above do not certify
that separate order-eight reduction.
The fresh restricted catalogs used below have $510$ ordinary classes
with nonzero rational $H_2$ at order eight and $8046$ at order nine;
their separate exhaustion proofs are
Propositions~\ref{CP2:A-positive-corners-eight} and
\ref{CP2:A-positive-corners-nine}. For comparison the full published
counts are
\[
\begin{array}{c|r|r}
\text{order}&\text{ordinary isomorphism}&\text{isomorphism and opposition}\\\hline
8&1,668,997&858,977\\
9&3,685,886,630&1,844,075,697.
\end{array}
\]
The ordinary counts are the opened OEIS record A058129; the opposition
counts are A058133 and the abstract of Distler--Kelsey at the authors'
university repository\cite{OEIS_A058129,OEIS_A058133,CP2-DistlerKelsey2009}.
Thus neither restricted catalog is advertised as a full catalog in its
order. Our reduction to nonzero rational second homology is part of
every use of these restricted lists.
\end{remark*}

\subsection{What the literal cup certificates establish}

\begin{proposition}[Proved: the certificate obstruction]
\label{CP2:A-cup-certificate}
Let $C$ be a finite connected category. Suppose normalized
nerve chains over $\mathbb F_2$ give $H_1=0$ and
$\dim H_2=1$. If a nonboundary two-cocycle $u$ admits a
three-cochain $v$ with $\delta v=u\smile u$, or if
$H_3(C;\mathbb F_2)\ne0$, then $BC$ is not homotopy equivalent
to $\CP^2$.
\end{proposition}
\begin{proof}
For chronological nonidentity arrows $[a_1|\cdots|a_j]$, the
normalized boundary deletes the first or last arrow or composes
an adjacent pair, with the usual alternating signs. A term
containing an identity after composition is degenerate and is
zero. These are exactly the face formulas of the nerve, so
their chain homology is that of $BC$.
The Alexander--Whitney formula on a four-simplex is
\[
(u\smile u)[a|b|c|d]=u[a|b]\,u[c|d].
\]
A literal equality $\delta v=u\smile u$ therefore says that
the square of the generator of $H^2(C;\mathbb F_2)$ is zero.
The corresponding generator of $H^2(\CP^2;\mathbb F_2)$ has
nonzero square. The other obstruction follows from the
vanishing of $H_3(\CP^2;\mathbb F_2)$. Both are invariant
under homotopy equivalence.
\end{proof}

In the new linear-profile search, the initial filter and the
cochain producer use separate bar calculations. The independent
checker then reconstructs each typed category from its literal
table, verifies every associativity equation, recomputes the
first two Betti numbers, checks that the supplied cocycle is
not a coboundary, and evaluates its square-primitive equation
on every normalized four-simplex. For a positive-third-homology
certificate it recomputes the necessary fourth-boundary rank.
No finite-degree calculation is used to assert vanishing of an
uncomputed homology tail.
